\section{Điều SHAP chưa giải quyết}
\label{sec:ch04-dieu-shap-chua-giai-quyet}

Bất định thống kê của phép tổng hợp không phải là điều duy nhất còn lại. Các
tiên đề làm rõ cách chia một chênh lệch đã được định nghĩa, nhưng không
chọn thay người dùng chênh lệch nào đáng chia. Giá trị nền, quy tắc hoàn thiện
đặc trưng còn thiếu và việc dùng kỳ vọng có điều kiện hay biên đều nằm trước
phương trình Shapley. Nếu những lựa chọn ấy không hợp với câu hỏi miền, một lời giải
thích vẫn có thể thỏa tiên đề mà không mô tả can thiệp người đọc quan tâm.

Sự phân biệt này đặc biệt quan trọng với \tn{giá trị Shapley}{Shapley value}
khi các đặc trưng phụ thuộc. Một đóng góp có thể phản ánh thông tin trùng lặp,
quan hệ tương quan hoặc một tình huống hoàn thiện chưa từng xuất hiện trong dữ
liệu. Các tiên đề bảo đảm \tn{tính nhất quán}{consistency} của phép chia trên hàm giá trị đã
chọn; chúng không xác nhận hàm giá trị ấy là mô hình nhân quả hoặc là thước đo
\tn{độ trung thực}{faithfulness}.

Vì vậy, SHAP là bước tiến so với việc chọn hàm nhân của LIME tùy ý, nhưng chưa
phải là điểm kết thúc của lập luận. Chương~\ref{ch:attribution-theo-gradient}
chuyển sang attribution theo gradient, nơi đường đi từ baseline tới đầu vào
thật lại đưa một lựa chọn nền khác vào phương trình. Qua các phương pháp, câu hỏi giữ
nguyên: thành phần vắng mặt đã được thay bằng điều gì, và vì sao sự thay thế đó
đại diện cho thay đổi có nghĩa?

\index{giá trị Shapley|)}%
